\section{The complete rank-three congruence integer}
\label{sec:rank-three-congruence-v46}

This chapter reconstructs the original $\mathrm{CNR3}$ from
\texttt{CongruenceNumberRankThree.lean}, including its two uniform
trace catalogues.  Its elementary product argument will choose the
integer before the coefficient prime.  The catalogue proof uses only
ranks three and two; it does not assume the original all-rank
$\mathrm{CountK}$ or $\mathrm{CountQ}$ producers.

Fix an imaginary quadratic field $K$.  Put $d=|\operatorname{disc}K|$.
Write $\varepsilon$ for the cyclotomic character, with
$\varepsilon(\sigma)=q$ at an arithmetic Frobenius $\sigma$ at a
place of residue cardinality $q$ away from its coefficient prime.
Consequently geometric Frobenius is $\sigma^{-1}$.
For a spherical automorphic representation of rank $n$, our
algebraic Satake multiset is
\[
 A_v(\pi)=q_v^{(n-1)/2}\operatorname{Sat}_v(\pi).
                                                               \tag{RNC.1}
\]
It is the geometric Frobenius multiset.  The arithmetic multiset
$\beta_v(\pi)$ is its elementwise inverse.

The original $\mathrm{Level}(\ell,M,m,r)$ is retained literally.
It includes continuity; away from $\ell$, inertia is diagonal with
characters killing the subgroup which fixes every $M$-th root of
unity.  At each place above $\ell$ there is a complete ordered
Galois flag with diagonal characters $\kappa_i$ and
\[
 \kappa_i|_I=\varepsilon^{m_i}\chi_i|_I,
 \qquad \chi_i\text{ kills the subgroup fixing }\mu_M,
                                                               \tag{RNC.2}
\]
together with the disjunction
\[
 (\forall i,j,\ \kappa_i\ne\kappa_j\varepsilon)
       \quad\text{or}\quad r\text{ is potentially crystalline}.
                                                               \tag{RNC.3}
\]
The potentially crystalline clause has its original local bridge
and period-ring meaning.  No restriction such as $\ell\mid M$ is
added.  All fields, algebraic closures, scalar towers and algebraic
isomorphisms $\iota:\overline{\mathbb Q}_\ell\simeq\mathbb C$
below are the original ones.  In particular $\ell=3$ is allowed.

We use the named classical foundations already recorded with CX:
Clozel's cohomology and rationality theorem, Borel--Serre finite
arithmetic CW models, Borel--Labesse--Schwermer's injection of
cuspidal cohomology, strong multiplicity one and local genericity,
local Langlands conductor and spherical compatibility, the generic
newvector theorem, and finite-coefficient $p$-adic Hodge comparison
and weak admissibility.  Their numerical and normalization
applications are proved below, rather than supplied as a counting
theorem.  The extra rank-two adelization step uses the cuspidal
decomposition in Jacquet--Langlands, Proposition 10.9, and the
explicit classical-to-adelic construction in RNC2.  The exact
published statements and primary URLs are recorded separately in
the classical input registry.

% BEGIN RNC1
\paragraph{RNC1. The precise automorphic Galois input in the two required ranks.}
\label{claim:RNC1}
For an initial regular algebraic cuspidal polarized pair over
a CM or totally real field, with arbitrary initial integral
highest weight $\lambda$, the published theorem supplies a
continuous semisimple Galois representation at every prime,
the away-from-that-prime comparison (RNC.5), and Hodge--Tate
weights $\{\lambda_i+n-i\}$.  This general existence statement
is available before a flag has identified $\lambda$.
After that identification, its following bounded
specialization will be used.
For the following two types of cuspidal representations, the
automorphic Galois representation exists for every coefficient
prime and every coefficient isomorphism, is continuous and
semisimple, and has the stated Hodge--Tate multiset:
\[
\begin{array}{c|c|c|c}
 F,n& m&\lambda_i=-m_{n+1-i}-(n-i)&\mathrm{HT}\text{ multiset}\\ \hline
 K,3&(e-1,-1,-e-1)&(e-1,0,1-e)&\{e+1,1,1-e\}\\
 \mathbb Q,2&(e,-e)&(e-1,-e)&\{e,-e\}.
\end{array}                                                     \tag{RNC.4}
\]
Here $e\geq1$, the first representation is conjugate self-dual,
and the second is essentially self-dual.  Their infinitesimal
character is $\operatorname{infOf}(m)=\{(n-1)/2+m_i\}$ at every
embedding.  Local compatibility, including monodromy, is
\[
 \iota\,\mathrm{WD}(R|_{G_{F_v}})^{\mathrm{F\!ss}}
      \simeq\operatorname{rec}_{F_v}
                    (\pi_v|\det|^{(1-n)/2}) .                 \tag{RNC.5}
\]
At $v\nmid\ell$ the parameter is pure, of geometric weight $2$
in the first row and $0$ in the second.  A spherical place above
the coefficient prime gives a crystalline representation.
\begin{proof}
This is a checked application of the published theorems, not a
new automorphy assertion.  Barnet-Lamb--Gee--Geraghty--Taylor,
\emph{Potential automorphy and change of weight}, Theorem 2.1.1,
applies to regular algebraic cuspidal polarized pairs over CM
\emph{or totally real} fields.  Its definition in Section 2.1
asks for $\pi^c\simeq\pi^\vee\otimes(\chi\circ N_{F/F^+})$,
with the prescribed constant infinity sign.  For the first row
take the quadratic character of $K/\mathbb Q$ as multiplier:
its pullback by the norm is trivial and its real sign is $-1$,
as required in odd rank.  For rank two over $\mathbb Q$ take
the actual central character whenever the cusp satisfies
the stated essential self-duality.  RNC2 will verify this
input for the newform application; it is not an antecedent
for the general existence theorem here.  There is only one real place, so the
constant-sign condition imposes no extra choice there.

The weight convention of that theorem is that $\pi_\infty$
has the infinitesimal character of the dual of the highest-weight
module of weight $\lambda$.  The formula in (RNC.4) gives exactly
$\lambda_i+n-i=-m_{n+1-i}$, so it matches the original
$\operatorname{infOf}$ convention.  Both displayed $\lambda$
are dominant integral for $e\geq1$.  The theorem's Hodge formula
and conjugate symmetry give respectively geometric weights $2$
and $0$.  Its clauses (2)--(4) give purity and (RNC.5) away from
the coefficient prime, and crystallinity at spherical places
above it.  These statements require neither irreducibility of
every associated representation nor an automorphy-lifting lower
bound on the prime.

For (RNC.5) at the coefficient prime use the same authors'
\emph{Local-global compatibility for $l=p$, II}, Theorem A,
Theorem 1.1 and Corollary 1.2.  They give full Frobenius
semisimplification for Shin-regular weight.  Their definition is:
odd rank suffices; in even rank a strict inequality
$\lambda_i>\lambda_{i+1}$ for some odd $i$ at some embedding
suffices.  Rank three meets the first alternative even at $e=1$.
For rank two, $\lambda_1-\lambda_2=2e-1>0$, so the second
alternative holds.  No Iwahori hypothesis remains in this
theorem, and every prime is allowed.  Its reciprocity sends a
uniformizer to geometric Frobenius, giving precisely (RNC.1)
and (RNC.5).  We use this stronger conclusion only in these
two verified ranks and weights.
\end{proof}
% END RNC1

% BEGIN RNC2
\paragraph{RNC2. The actual newform supplies the required rank-two cusp.}
\label{claim:RNC2}
Let $e\geq1$, let $f$ be the original newform of weight
$k=2e+1$, and suppose $f$ is weakly attached to
$W\varepsilon^{-e}$ in the original sense.  Then there is a
regular algebraic essentially self-dual cusp $\pi$ on
$\mathrm{GL}_2/\mathbb Q$, of the second infinity type in
(RNC.4), weakly attached to the same $W$.  At every prime
$p\nmid f.M$, if $\alpha_p,\alpha'_p$ are the roots of
$f.\mathrm{heckePoly}(p)$, its algebraic Satake multiset is
\[
           A_p(\pi)=\{p^{-e}\alpha_p,p^{-e}\alpha'_p\}.
                                                               \tag{RNC.6}
\]
\begin{proof}
Here is the bounded adelization and its normalization.  The
original definition supplies the $\Gamma_0(f.M)$ nebentype law,
cuspidality, $a_1(f)=1$, and the coefficient identity
\[
 a_{pn}(f)+\chi(p)p^{k-1}{\bf1}_{p\mid n}a_{n/p}(f)
                         =a_p(f)a_n(f)\quad(p\nmid f.M).
                                                               \tag{RNC.7}
\]
On $g_\infty\in\mathrm{GL}_2^+(\mathbb R)$ put
\[
 F_\infty(g_\infty)=\det(g_\infty)^{k/2}
               j(g_\infty,i)^{-k}f(g_\infty i),
 \qquad j\!\left(\begin{pmatrix}a&b\\c&d\end{pmatrix},z\right)
                     =cz+d.
                                                               \tag{RNC.8}
\]
Strong approximation writes every adelic matrix as
$\gamma g_\infty u$, with $\gamma\in\mathrm{GL}_2(\mathbb Q)$
and $u\in K_0(f.M)$.  Define its value to be
$\chi(d_u)^{-1}F_\infty(g_\infty)$.  Two such decompositions
differ by an integral determinant-one matrix in
$\Gamma_0(f.M)$; its nebentype factor in (RNC.8) cancels the
change in $\chi(d_u)^{-1}$.  Thus this is a well-defined actual
left-rational-invariant adelic cusp form.  Its compact-open
finite transformation law is explicit.  Its unipotent constant
terms are the classical cusp constant terms and vanish.
The factor in (RNC.8) is invariant under positive real scalar
matrices and has rotation weight $k$; differentiating the
holomorphic identity gives the usual lowest discrete-series
vector, with infinitesimal character
$\{(k-1)/2,-(k-1)/2\}$.

Decompose this nonzero form in the cuspidal representation space
with its central character, by Jacquet--Langlands,
Proposition 10.9.  The fixed finite level and fixed infinity
type make the relevant space finite-dimensional, so at least
one irreducible cuspidal summand has a nonzero projection of
the form.  At $p\nmid f.M$, the $p+1$ standard double-coset
representatives act on (RNC.8) by the two terms of (RNC.7).
The normalized spherical parameters of every such nonzero
projection therefore have sum $p^{-(k-1)/2}a_p(f)$ and product
$\chi(p)$.  This also follows directly by comparing every
Fourier coefficient in (RNC.7); a holomorphic cusp form with
all coefficients zero vanishes.  These identities are preserved
by the projection, since the spherical Hecke algebra commutes
with the global representation decomposition.

Let $\pi^0$ be one of these summands and set
$\pi=\pi^0|\det|^{1/2}$.  At the finite uniformizer this twist
multiplies each spherical root by $p^{-1/2}$, while at infinity
it adds $1/2$ to both infinitesimal parameters.  As $k=2e+1$,
the infinity multiset becomes $\{e+1/2,1/2-e\}$ and (RNC.1)
gives (RNC.6).  This is the exact original
$\mathrm{NewformRep}$ specialization $j=e$, not a change to
the original form or Galois representation.

For any rank-two unramified parameter with product equal to
the central-character value, inversion is the same unordered
pair as multiplication by the inverse central character.
Strong multiplicity one consequently gives
$\pi^\vee\simeq\pi\otimes\omega_\pi^{-1}$.
The central character is algebraic at infinity for this
displayed integral highest weight, so $\pi$ is an admissible
polarized pair for RNC1.  Finally, at geometric Frobenius
$\varepsilon^{-e}$ has value $p^e$.  The original weak
attachment of $f$ therefore says that the eigenvalues of
$W(\sigma^{-1})$ are exactly (RNC.6), at the same cofinite
set of primes and with all original prolongations and lifts.
This proves the required weak attachment to $W$.
\end{proof}
% END RNC2

% BEGIN RNC3
\paragraph{RNC3. Weak attachment identifies the actual representation and its weights.}
\label{claim:RNC3}
Suppose an original continuous irreducible representation $r$
is weakly attached to an initial regular algebraic
conjugate-self-dual cusp in rank three, or to the rank-two
cusp of RNC2, and has the corresponding original level
and ordered weight $m$ from (RNC.4).  Then it
is isomorphic, in its original algebraic closure, to the
representation supplied by RNC1.  Its original complete flag
recovers exactly the highest weight in (RNC.4).  The statement
applies in rank three to $r=\rho\varepsilon^{-1}$ from
$\mathrm{Adm3}(\iota,e,N,\rho)$, and in rank two to the
actual $W$ from $\mathrm{Adm2}$.
\begin{proof}
First choose the representation attached by the published
theorem to the given regular algebraic cusp; at this point its
highest weight need not have been identified: this uses
the general-existence part of RNC1, not its later Shin-regular
specialization.  Its cofinite
Frobenius characteristic polynomials agree with those of $r$
by the definition of weak attachment and (RNC.1).
Both representations are continuous and semisimple.  ADK1
captures their compact images in a common finite coefficient
extension.  Chebotarev and continuity make their traces equal
on the whole original Galois group.  The semisimple character
argument of CAP6 then gives an actual intertwining matrix over
$\overline{\mathbb Q}_\ell$.  This argument neither changes
$\iota$ nor replaces the chosen algebraic closure.

The representation is therefore de Rham.  Its flag subobjects
and successive quotients are de Rham by exactness of the
de Rham category.  Each quotient in (RNC.2) is an unramified
character times a finite character times
$\varepsilon^{m_i}$.  Its Hodge--Tate weight at every relevant
embedding is $-m_i$.  Exactness gives the whole multiset
$\{-m_i\}$; it is not inferred merely from being an extension
of de Rham characters.  The automorphic Hodge formula now
identifies the decreasing integers $\lambda_i+n-i$ with
$-m_{n+1-i}$, giving (RNC.4) and the asserted infinity type.
Every local comparison is transported by this same matrix.

For rank three, twisting by $\varepsilon^{-1}$ preserves
irreducibility and continuity.  Away from $\ell$ it changes
no inertia character.  At $\ell$ it replaces $m=(e,0,-e)$
by $(e-1,-1,-e-1)$ and every $\kappa_i$ by
$\kappa_i\varepsilon^{-1}$.  The first alternative in
(RNC.3) is unchanged; the second is preserved because the
cyclotomic character is crystalline.  Thus the same original
integer $N$ is a level for $r$, and its automorphy is exactly
the $\mathrm{Adm3.aut}$ premise.  For rank two, RNC2 gives
the required weak cusp for the original $W$; its original
level and irreducibility are fields of $\mathrm{Adm2}$.
\end{proof}
% END RNC3

% BEGIN RNC4
\paragraph{RNC4. The flag detects every forbidden monodromy ratio.}
\label{claim:RNC4}
At a place $v\mid\ell$ of residue cardinality $q$, let $r$
be as in RNC3.  The semisimple Weil characters of its parameter
are the rank-one parameters $\mu_i$ of the ordered flag.
With arithmetic Weil degree $\deg(s)$, they satisfy
\[
 \mu_i(s)=q^{m_i\deg(s)}\iota(\psi_i(s)),\qquad
 \psi_i=\kappa_i\varepsilon^{-m_i}.                            \tag{RNC.9}
\]
Their inertia restrictions are $\iota\chi_i$.  If
$\mu_i(s)=q^{\deg(s)}\mu_j(s)$ for all Weil elements $s$,
then $\kappa_i=\kappa_j\varepsilon$ on the whole local Galois
group.  Consequently the first alternative of (RNC.3) forces
zero monodromy, both for the Galois and the automorphic
parameters.
\begin{proof}
Use a common finite coefficient extension as in ADK1.
RNC3 makes $r$ and all its flag terms de Rham, hence potentially
semistable.  Exactness of potentially semistable comparison
takes the flag to a flag of filtered modules and of raw Weil
actions; on its successive rank-one quotients the cyclotomic
character has arithmetic Weil value $q^{\deg(s)}$, a finite
character retains its inertia value, and an unramified
character retains its Frobenius value.  This gives (RNC.9),
including every original local bridge after scalar extension.
Frobenius semisimplification preserves the characteristic
polynomial on every Weil element, so the semisimple Weil
representation is the direct sum of these characters.

The continuous character $\psi_i$ has finite inertia image:
its restriction there is $\chi_i$.  More generally a continuous
character of the compact local Galois group has norm one,
because its norm image is a compact subgroup of
$\mathbb R_{>0}$ and all integer powers of any element are
bounded.  A proposed ratio in the assertion would give
\[
 \iota(\psi_i/\psi_j)(s)
                    =q^{(1+m_j-m_i)\deg(s)}.
                                                               \tag{RNC.10}
\]
Take an arithmetic Frobenius lift and apply $\iota^{-1}$,
which fixes the rational integer $q$.  The left side has
$\ell$-adic norm one, whereas $q$ has norm strictly less than
one.  Thus $m_i=m_j+1$.  Equation (RNC.10) then says that
$\psi_i/\psi_j$ is one on the Weil group.  Density and
continuity make it one on the whole decomposition group,
giving the claimed equality of $\kappa$'s.

In a Frobenius-semisimple Weil--Deligne module, a nonzero
monodromy entry maps some character space into its norm twist,
by $r(s)Nr(s)^{-1}=q^{\pm\deg(s)}N$; interchanging the two
indices covers the sign convention.  All such ratios were
excluded.  Therefore $N=0$.  Equality of the semisimple Weil
representations in (RNC.5) gives exactly the same exclusion
for the automorphic parameter.  The argument does not assert
that arbitrary extensions of characters are crystalline.
\end{proof}
% END RNC4

% BEGIN RNC5
\paragraph{RNC5. A polynomial automorphic level independent of the coefficient prime.}
\label{claim:RNC5}
The same cusps associated to the original representations have
nonzero principal-level fixed vectors at $M^{2n-1}$: level
$N^5$ for $\mathrm{Adm3}$ and level $M^3$ for the $W$ in
$\mathrm{Adm2}(\iota,e,M,\rho)$.  They are spherical at
every place not dividing $M$.  The original Galois level
remains $M$ throughout.
\begin{proof}
For a finite place $v$ write $a_v=v(M)$, normalized so that
$M\mathcal O_{F_v}=\mathfrak p_v^{a_v}$.  A character which
kills the subgroup fixing $\mu_M$ has local conductor at most
$a_v$.  Indeed local reciprocity identifies its restriction
with a character through the cyclotomic extension.  For
$x\in1+M\mathcal O_{F_v}$, the norm to the rational local
field is congruent to $1$ modulo $M$: multiplication by $x$
on an integral local basis has matrix $1+M B$, whose
determinant is $1$ modulo $M$.  The cyclotomic action is
therefore trivial.  If $a_v=0$, inertia fixes the prime-to-
residue-characteristic roots in question, so the conductor
is zero.  This proof also covers ramified places of $F$.

At $v\nmid\ell$ the original inertial action is a direct sum
of these finite characters.  Its monodromy is zero and its
Artin conductor is at most $n a_v$.  At $v\mid\ell$, RNC4
identifies the semisimple inertial characters of the parameter
with the same finite characters, so their conductor sum is
again at most $n a_v$.  For an arbitrary nilpotent monodromy
operator the full conductor is
\[
 a(r_v,N_v)=a(r_v)+\dim V^{I_v}-\dim(\ker N_v)^{I_v}
                              \leq n a_v+n-1.                 \tag{RNC.11}
\]
If $a_v\geq1$ this is at most $(2n-1)a_v$.

It remains to exclude an unwanted coefficient-prime factor
when $a_v=0$.  Under the first alternative of (RNC.3), RNC4
gives $N_v=0$.  Under the second, the Galois parameter has
zero monodromy by potential crystallinity, and the full
compatibility in RNC1 transfers this to the automorphic
parameter.  Inertia is already trivial when $a_v=0$.
The parameter is unramified, hence the local cusp component
is spherical.  This proves conductor at most $(2n-1)a_v$
at every place, with no condition $\ell\mid M$.

Every local component of a cuspidal $\mathrm{GL}_n$
representation is generic.  Local Langlands preserves its
conductor, and the generic newvector theorem supplies a
nonzero vector fixed by $K_1(\mathfrak p_v^{b_v})$ at its
conductor $b_v$.  The principal subgroup of depth
$(2n-1)a_v$ is contained in this subgroup.  Tensor these
vectors over the finitely many nonspherical places, using
the spherical vectors elsewhere.  The restricted tensor
product realization of the actual cusp gives a nonzero
global vector at the asserted principal level.  Thus this
is a bound on the auxiliary automorphic label set, not a
replacement of the original Galois level condition.
\end{proof}
% END RNC5

% BEGIN RNC6
\paragraph{RNC6. Fixed cells and explicit coefficient ranks give uniform finite catalogues.}
\label{claim:RNC6}
For fixed $K$ there are natural constants $c_3,c_2$ such that,
for every $e\geq1$ and $L\geq1$, the numbers of cusp
isomorphism classes with principal-level $L$ vectors and the
two infinity types of (RNC.4) are at most
\[
 \begin{split}
 c_3 L^{18}(e+1)^{18}&\quad(F=K,n=3),\\
 c_2 L^4(e+1)^4&\quad(F=\mathbb Q,n=2).
 \end{split}                                                   \tag{RNC.12}
\]
One may restrict the first set to conjugate-self-dual classes.
\begin{proof}
We give the bounded-rank version of the fixed-cell argument
underlying CX1.  In the first row, the highest-weight module
of weight $(e-1,0,1-e)$ is a determinant twist of the module
of weight $(2e-2,e-1,0)$.  HP1--HP2 realize the latter in
polynomials of degree $3e-3$ in the nine matrix coordinates.
Its dimension is at most $(3e-2)^9\leq[3(e+1)]^9$; this
inequality follows even by bounding each individual monomial
exponent between zero and $3e-3$.  Tensor over the two
embeddings of $K$ to obtain coefficient rank at most
$[3(e+1)]^{18}$.  The determinant twist changes no rank.
For the second row, shift $(e-1,-e)$ by $e$ to obtain
$(2e-1,0)$.  The same polynomial realization in four matrix
coordinates bounds the coefficient rank by
$(2e)^4\leq[2(e+1)]^4$.  These formulas include $e=1$.

Use principal level $3$ once and for all.  Its arithmetic
components may stabilize projective $\mathcal O_F$-lattices
rather than the standard free one.  Each is nevertheless
a free integral lattice of rank $dn$ after restriction of
scalars to $\mathbb Z$, and its principal level-$3$ action
is congruent to one on that same lattice.  Thus its arithmetic
groups are torsion-free: restriction of scalars embeds them
in the kernel of $\mathrm{GL}_{dn}(\mathbb Z)\to
\mathrm{GL}_{dn}(\mathbb Z/3)$, with $d=[F:\mathbb Q]$.
For completeness, if a finite-order matrix in this kernel
is not one, take a prime-order power of it and write it as
$1+3^a B$, $a\geq1$, with an entry of $B$ prime to $3$.
Expanding its prime-th power, for a prime other than $3$
the linear term has strictly smaller $3$-adic valuation
than all later terms; for prime $3$ the same holds because
$a\geq1$.  It cannot equal one.  This is a contradiction.

The level-$3$ adelic quotient has finitely many components,
using the finite ideal class group.  Remove the rational
split central direction.  Borel--Serre gives a finite CW
model for each of these fixed torsion-free components; let
$b_F$ be their total number of cells, in all degrees.
The cover at level $3L$ has total covering degree at most
\[
 [K_f(3):K_f(3L)]\leq
 |M_n(\mathcal O_F/(3L))|=(3L)^{dn^2}.                         \tag{RNC.13}
\]
The sum of the ranks of cellular cochains with coefficient
rank $D$ is therefore at most $b_F(3L)^{dn^2}D$.
Their total cohomology dimension is no larger.

Every cusp under consideration is cohomological for the
explicit coefficient just constructed, by Clozel's theorem
with the matching infinitesimal character.  A nonzero
level-$L$ vector is also a level-$3L$ vector.  The cuspidal
cohomology injection of Borel--Labesse--Schwermer places its
nonzero cohomological summand in ordinary cohomology at
that level.  Distinct global cusp classes give independent
summands; multiplicity one, or simply decomposition into
inequivalent irreducible cuspidal summands, ensures that
each class costs at least one dimension.  Applying (RNC.13)
with $dn^2=18$ and $4$, and the coefficient-rank bounds,
proves (RNC.12) after enlarging two fixed integer constants.
Neither constant depends on the coefficient prime, $e$, $L$,
or the place at which a Frobenius value will later be read.
\end{proof}
% END RNC6

% BEGIN RNC7
\paragraph{RNC7. Full cusp labels and algebraic Satake conjugation.}
\label{claim:RNC7}
Let $S_3(e,N)$ be the set of conjugate-self-dual cusp classes
of the first type in (RNC.4) with level-$N^5$ vectors.
Let $S_2(e,M)$ be all cusp classes of the second type with
level-$M^3$ vectors.  They are finite and are permuted by
every automorphism of $\mathbb C$.  At a place prime to the
respective level, this permutation carries $A_v$ and
$\beta_v$ to their coefficientwise conjugates.  Moreover
\[
 |S_3(e,N)|\leq c_3[N(e+1)]^{90},\qquad
 |S_2(e,M)|\leq c_2[M(e+1)]^{12}.                              \tag{RNC.14}
\]
\begin{proof}
Finiteness and the inequalities follow from RNC6, since
$N^{90}(e+1)^{18}\leq[N(e+1)]^{90}$ and
$M^{12}(e+1)^4\leq[M(e+1)]^{12}$.
The labels are full global cusp classes, not eigenvalues
at a single prime.  Clozel's rationality and conjugation
theorem applies to these regular $C$-algebraic cusps with
the displayed matching coefficient.  Both embeddings of
$K$ have the same weight, and $\mathbb Q$ has one embedding,
so the infinity condition survives every field automorphism.
The principal fixed-vector condition survives scalar
conjugation as well.  Duality, conjugation of $K$, and
scalar conjugation commute on the finite representation;
strong multiplicity one then preserves conjugate
self-duality in the first catalogue.  In rank two every
class is essentially self-dual by the last argument of
RNC2.  Applying the inverse field automorphism gives the
inverse permutation of each finite catalogue.

There is a possible square-root trap in raw normalized
Satake parameters, so we record the exact calculation.
If $T_{v,j}$ is the unnormalized characteristic-function
Hecke operator for
$\operatorname{diag}(\varpi^{[j]},1^{[n-j]})$ at spherical
level, then
\[
       e_j(A_v)=q_v^{j(j-1)/2}\,T_{v,j}(\pi).                 \tag{RNC.15}
\]
Indeed normalized Satake gives
$T_{v,j}=q_v^{j(n-j)/2}e_j(\operatorname{Sat}_v)$;
substitution in (RNC.1) proves (RNC.15).  The exponent in
(RNC.15) is an integer.  The Hecke correspondence has
rational coefficients, so scalar conjugation applies to
its eigenvalue and fixes that rational power of $q_v$.
The monic characteristic polynomial consequently has
coefficientwise conjugate roots.  Inverting them proves
the same assertion for $\beta_v$.  This also proves that
all these roots are algebraic, since the finite parts have
number-field models.  No automorphism is required to fix
an arbitrarily chosen square root of $q_v$.
\end{proof}
% END RNC7

% BEGIN RNC8
\paragraph{RNC8. Purity and both integral endpoint bounds.}
\label{claim:RNC8}
For the first catalogue at a place $w\mid p$ with
$p\nmid N$, each arithmetic root $b\in\beta_w$ has
$|b|=q^{-1}$ and both
\[
          q^{e+1}b,\qquad q^{e-1}b^{-1}
                       \quad\text{algebraic integers}.       \tag{RNC.16}
\]
For the second catalogue at $p\nmid M$, each arithmetic
root $b\in\beta_p$ has $|b|=1$ and both
\[
                    p^e b,\qquad p^e b^{-1}
                       \quad\text{algebraic integers}.       \tag{RNC.17}
\]
These assertions hold after every complex field automorphism.
\begin{proof}
The indicated places are spherical because principal level
is maximal compact there.  Apply RNC1 at a coefficient
prime different from $p$.  The parameter has zero monodromy,
so purity gives absolute value $q$ for geometric roots in
rank three, and $1$ in rank two.  Inversion gives the
displayed values for $b$.  RNC7 gives the same argument for
every conjugate label and hence every algebraic conjugate.

To check integrality, fix an embedding of the algebraic
roots into $\overline{\mathbb Q}_a$ for each rational prime
$a$.  For $a\ne p$ use the representation of RNC1 at
coefficient prime $a$.  After ADK1 captures its compact
image in a finite extension, the image preserves an
integral lattice.  One elementary construction is to take
the module spanned by the translates of a starting lattice:
compactness bounds them in one larger lattice, which is
Noetherian, so this span is a lattice.  Frobenius and its
inverse preserve it.  Their eigenvalues are thus integral
units, proving all required integrality away from $p$.

At $a=p$, use the same cusp's crystalline representation
from RNC1.  For a weakly admissible filtered Frobenius
module whose Hodge--Tate weights lie in $[h_-,h_+]$, all
Newton slopes lie in that interval.  Here is the needed
argument: on the minimal-slope subobject,
$t_N\geq t_H\geq h_-\dim$; the maximal-slope quotient has
$t_N\leq t_H\leq h_+\dim$, by the subobject inequality
and equality for the whole module.  Passing to a splitting
coefficient field and to a residue-degree power of Frobenius
gives the same inequalities for each root, with valuation
normalized by $v(q)=1$.

In rank three the geometric interval is $[1-e,e+1]$;
the arithmetic interval after inversion is
$[-e-1,e-1]$.  This proves (RNC.16) at every prime above
$p$.  In rank two both intervals are $[-e,e]$, giving
(RNC.17).  Together with the preceding unit calculation
these bounds prove integrality at every finite place of
a number field containing the roots, which is equivalent
to algebraic integrality.  The argument works for every
embedding and every conjugate catalogue label.  In
particular it proves all the original $\mathrm{EigData.int}$
and $\mathrm{intInv}$ clauses: their permitted exponents
are respectively at least $e+1,e-1$ in rank three and
at least $e,e$ in rank two.
\end{proof}
% END RNC8

% BEGIN RNC9
\paragraph{RNC9. The rank-three trace catalogue classifies every later coefficient realization.}
\label{claim:RNC9}
For fixed $e,N,p,w$ as above let $q=\operatorname{N}w$.
The already chosen labels $S_3(e,N)$ carry traces
\[
                    t_3(s)=\sum_{b\in\beta_w(s)}(qb)^{-1}.
                                                               \tag{RNC.18}
\]
They satisfy $|t_3(s)|\leq3$, $q^e t_3(s)$ is an algebraic
integer, and every automorphism of $\mathbb C$ permutes
the traces with their labels.  For every later odd
$\ell\ne p$, closure, $\iota$, prolongation $B$ of $w$
and arithmetic Frobenius $\sigma$, every
$\mathrm{Adm3}(\iota,e,N,\rho)$ has a label with
\[
        \operatorname{tr}\rho(\sigma^{-1})=\iota^{-1}t_3(s).
                                                               \tag{RNC.19}
\]
Equal labels for two such representations imply that they
are isomorphic.
\begin{proof}
RNC8 gives three summands of absolute value one in
(RNC.18).  It also shows that
$q^e(qb)^{-1}=q^{e-1}b^{-1}$ is integral.  RNC7 gives
the exact permutation assertion because $q$ is rational.
RNC3 identifies $r=\rho\varepsilon^{-1}$ with the
representation attached to its actual cusp, and RNC5
puts that cusp into the already fixed label set.
At $w\nmid\ell$ full local compatibility gives
unramifiedness at every prolongation and the geometric
multiset $\beta_w(s)^{-1}$.  Since
\[
             r(\sigma^{-1})=q\rho(\sigma^{-1}),
                                                               \tag{RNC.20}
\]
the trace is exactly (RNC.19), for every permitted lift.
If the two cusp labels agree, both twisted representations
are isomorphic to the same attached semisimple
representation in the same coefficient closure; equivalently
one may apply CAP6 directly to their cofinite polynomials.
Untwisting proves isomorphism of the original $\rho$'s.
Thus the map classifies full representations, even when
distinct labels happen to have the same trace at $w$.
\end{proof}
% END RNC9

% BEGIN RNC10
\paragraph{RNC10. The rank-two catalogue gives the exact restricted block trace.}
\label{claim:RNC10}
Let $M\geq1$, $p\nmid M$, and $q=p^h$ with $h\geq1$.
The labels $S_2(e,M)$ carry traces
\[
                  t_2(s)=1+\sum_{b\in\beta_p(s)}b^{-h}.
                                                               \tag{RNC.21}
\]
They have absolute value at most $3$, $q^e t_2(s)$ is
integral, and all complex automorphisms permute them with
multiplicity.  For every later original coefficient datum,
every $\mathrm{Adm2}(\iota,e,M,\rho)$ and every original
arithmetic Frobenius $\sigma$ at a $K$-place of norm $q$,
there is a classifying label with
\[
         \operatorname{tr}\rho(\sigma^{-1})=\iota^{-1}t_2(s).
                                                               \tag{RNC.22}
\]
Equal labels imply isomorphism of the rank-three block
representations, including their original ordering of blocks.
\begin{proof}
By RNC8 each of the two $b$'s has modulus one, so (RNC.21)
has modulus at most $3$.  Its scaled summands are
\[
 q^e t_2(s)=p^{he}+\sum_{b\in\beta_p(s)}(p^e b^{-1})^h,
                                                               \tag{RNC.23}
\]
which are integral.  RNC7 gives the permutation property.

Choose the actual $W$ and the actual conjugating matrix in
the original $\mathrm{Adm2}$ witness.  RNC2, RNC3 and RNC5
give its label in $S_2(e,M)$.  As $p\nmid M$ and
$p\ne\ell$, $W$ is unramified at $p$ and its geometric
Frobenius roots are $\beta_p(s)^{-1}$.  The restriction of
an arithmetic Frobenius of a residue-degree-$h$ extension
to $G_{\mathbb Q}$ differs from the $h$-th power of a
rational arithmetic Frobenius by inertia.  Applying $W$
therefore gives precisely the $h$-th powers of the same
two geometric roots.  In the ordered source matrix the
two-dimensional block occupies positions $0,2$, and the
middle diagonal entry is $1$.  Its trace is exactly
(RNC.21); conjugation by the given matrix preserves trace.

If labels agree, CAP6 or the common attached representation
identifies the two $W$'s.  Restrict its actual intertwiner
to $G_K$ and insert it in positions $0,2$, with $1$ in the
middle position.  Conjugating by the two original block
identifications gives an intertwiner of the two $\rho$'s.
No irreducibility of $W|_{G_K}$ is needed, and no converse
assertion about isomorphic restrictions is used.
\end{proof}
% END RNC10

% BEGIN RNC11
\paragraph{RNC11. One finite trace catalogue handles the original disjunction.}
\label{claim:RNC11}
For $e,N\geq1$ and a prime $p$ with $p\nmid Nd$, choose
$w\mid p$, $q=\operatorname{N}w$.  Before choosing any
coefficient prime, there is a finite set $S$ and numbers
$t_s\in\mathbb C$ such that
\[
 \begin{gathered}
 |S|\leq c_3[N(e+1)]^{90}
                   +c_2[Nd(e+1)]^{12},\\
 |t_s|\leq3,\qquad q^e t_s\text{ is an algebraic integer},
 \qquad \forall\tau\in\operatorname{Aut}(\mathbb C),\
          \exists u_\tau:S\simeq S,\quad
          \tau(t_s)=t_{u_\tau(s)}.
 \end{gathered}                                                 \tag{RNC.24}
\]
For every later original coefficient datum and Frobenius,
the predicate
$\mathrm{Adm3}(\iota,e,N,\rho)\lor
\mathrm{Adm2}(\iota,e,Nd,\rho)$ is classified by these
labels with trace $\iota^{-1}t_s$ and equal-label
isomorphism.
\begin{proof}
The residue degree of $w$ is $h\in\{1,2\}$ because $K$
is quadratic, so $q=p^h$.  Take the disjoint union of the
rank-three catalogue of RNC9 and the rank-two catalogue
of RNC10 at $M=Nd$.  Both are available since $p\nmid Nd$.
All three arithmetic assertions in (RNC.24) and its
cardinality bound follow summand by summand.  For a
representation satisfying the disjunction, use its
rank-three classifier when the first predicate holds;
otherwise use a chosen rank-two witness and its
classifier.  Equal labels cannot lie in different
summands, and the corresponding individual classifier
gives isomorphism.  This convention also covers a
representation satisfying both predicates, without
assuming the two classes of representations are disjoint.
The sets and their traces were constructed entirely from
$K,e,N,p,w$, before $\ell$, $\Omega$ or $\iota$.
\end{proof}
% END RNC11

% BEGIN RNC12
\paragraph{RNC12. Every Eisenstein comparison factor is nonzero and integral.}
\label{claim:RNC12}
Let $q\geq7$, $e\geq1$, let $z$ be a root of unity,
and suppose $|t|\leq3$ and $q^e t$ is an algebraic integer.
For the exact original expression
\[
 E(q,e,z)=q^{-e}z+1+q^e z^{-1},\qquad
              F=q^e\bigl(t-E(q,e,z)\bigr),                    \tag{RNC.25}
\]
$F$ is a nonzero algebraic integer and
\[
                         \log|F|\leq2e\log q+1.
                                                               \tag{RNC.26}
\]
\begin{proof}
A root of unity has modulus one, and its inverse is also
an algebraic integer.  The reverse triangle inequality
gives
$|E(q,e,z)|\geq q^e-1-q^{-e}\geq7-1-1/7>3$.
Therefore $t\ne E(q,e,z)$ and $F\ne0$.
Expand
\[
             F=q^e t-z-q^e-q^{2e}z^{-1}.                      \tag{RNC.27}
\]
Every term is integral.  If $T=q^e\geq7$, the triangle
inequality gives
$|F|\leq T^2+4T+1\leq2T^2$, since
$1+4/7+1/49<2$.  Taking logarithms, using
$\log2\leq1$ and $\log T=e\log q$, proves (RNC.26).
This proof needs no lower bound on individual trace gaps
at a coefficient prime.
\end{proof}
% END RNC12

% BEGIN RNC13
\paragraph{RNC13. The whole orbit product is a nonzero rational integer.}
\label{claim:RNC13}
Let $(S,t)$ be the catalogue in (RNC.24), and let
$z_o$ be roots of unity indexed by any finite set $\mathcal O$,
with the original condition that every automorphism of
$\mathbb C$ permutes this indexed family, with multiplicity.
Then
\[
       NN=\prod_{o\in\mathcal O}\prod_{s\in S}
                   q^e\bigl(t_s-E(q,e,z_o)\bigr)\in\mathbb Z
                                                               \tag{RNC.28}
\]
is nonzero and satisfies
\[
      \log|NN|\leq |\mathcal O|\,|S|\,(2e\log q+1).
                                                               \tag{RNC.29}
\]
\begin{proof}
Every factor is a nonzero algebraic integer by RNC12, so
the finite product is a nonzero algebraic integer.  For
an arbitrary field automorphism, use its permutation of
$\mathcal O$ and the independent permutation of $S$ in
(RNC.24).  Since $q$ is rational and the formula uses only
integer powers and inverses, these permutations carry
the factors of (RNC.28) bijectively to their conjugates.
The whole product is fixed.  An algebraic complex number
fixed by every complex field automorphism is rational:
any nontrivial algebraic conjugate extends to such an
automorphism by extending a transcendence basis.  A
rational algebraic integer is in $\mathbb Z$.
Finally the logarithm of the norm of a finite product
of nonzero complex numbers is the sum of the logarithms;
RNC12 gives (RNC.29).  Empty index sets cause no exception:
the empty product is $1$ and the logarithmic bound is $0$.
No coefficient prime has yet been chosen.
\end{proof}
% END RNC13

% BEGIN RNC14
\paragraph{RNC14. Selecting pairwise nonisomorphic families gives the required local product inequality.}
\label{claim:RNC14}
Choose any of the later original data with an odd
$\ell\ne p$, and for each $o$ let $X_o$ be a finite
family of the representations classified by RNC11,
pairwise nonisomorphic within that family.  Then every
original trace difference is nonzero and
\[
 |NN|_\ell\leq
 \prod_{o\in\mathcal O}\prod_{\rho\in X_o}
 \left|\operatorname{tr}\rho(\sigma^{-1})-
                         \iota^{-1}E(q,e,z_o)\right|_\ell.
                                                               \tag{RNC.30}
\]
\begin{proof}
For this fixed coefficient datum and Frobenius, choose
the classifier from RNC11.  Its restriction to each
$X_o$ is injective: equal labels imply isomorphism, and
the hypothesis on $X_o$ then identifies the two members.
The map $(o,\rho)\mapsto(o,s_\rho)$ is therefore an
injection into $\mathcal O\times S$.  Repeating the same
representation for different $o$ is harmless because
the first coordinate keeps the chosen factors distinct.
RNC12 and the exact trace identity in the classifier
give nonvanishing of every displayed difference.

Apply $\iota^{-1}$ to (RNC.28).  A monic integral
polynomial and the ultrametric inequality imply that
every algebraic integer in $\overline{\mathbb Q}_\ell$
has norm at most one.  Thus each factor has norm at most
one.  Removing the factors outside the indicated
injection can only increase the product of norms.
Finally $q=p^h$ and $\ell\ne p$, so $|q^e|_\ell=1$.
The norm of each selected factor is precisely the
corresponding unscaled trace-gap norm in (RNC.30).
This proves the inequality for every family at once,
with the same already chosen integer $NN$.
\end{proof}
% END RNC14

% BEGIN RNC15
\paragraph{RNC15. The constants are chosen before all varying arithmetic data.}
\label{claim:RNC15}
There are natural constants $c,A$, depending only on $K$,
such that (RNC.29) is at most
\[
                c\,|\mathcal O|\,[N(e+1)]^A\log p.
                                                               \tag{RNC.31}
\]
One permissible choice is
\[
             A=103,\qquad c=5(c_3+c_2d^{12}),                 \tag{RNC.32}
\]
with the fixed constants of RNC6 enlarged to natural
integers if necessary.
\begin{proof}
Put $Y=N(e+1)\geq1$.  Formula (RNC.24) gives
\[
 |S|\leq c_3Y^{90}+c_2d^{12}Y^{12}
                    \leq(c_3+c_2d^{12})Y^{102}.               \tag{RNC.33}
\]
Here the deliberately nonoptimal exponent is the sum
of the two catalogue exponents.  Since $q=p^h$ with
$1\leq h\leq2$, $\log q\leq2\log p$.  For $p\geq7$,
$\log p\geq1$, and $e+1\leq Y$.  Consequently
\[
 2e\log q+1\leq4e\log p+\log p
                         \leq5Y\log p.                       \tag{RNC.34}
\]
Insert (RNC.33)--(RNC.34) into (RNC.29).  The degree,
discriminant and cell constants were fixed with $K$;
none of $e,N,p,w,\mathcal O,z$, or any later coefficient
data occurs in (RNC.32).
\end{proof}
% END RNC15

% BEGIN RNC16
\paragraph{RNC16. The literal original $\mathrm{CNR3}$.}
\label{claim:RNC16}
For every original imaginary quadratic field $K$, choose
$c,A$ as in RNC15.  For every $e,N\geq1$, every rational
prime $p\geq7$ with $p\nmid N$ and
$p\nmid\operatorname{disc}K$, every $w\mid p$, and
every original finite root-of-unity orbit with its
full complex-automorphism permutation condition, there
exists $NN\in\mathbb Z\setminus\{0\}$ satisfying
(RNC.31).  For every subsequent odd $\ell\ne p$,
every original coefficient closure and scalar tower,
every $\iota$, every prolongation of $w$, every
arithmetic Frobenius $\sigma$, and every family
$X_o$ of pairwise nonisomorphic representations satisfying
\[
 \mathrm{Adm3}(\iota,e,N,\rho)
          \quad\text{or}\quad
 \mathrm{Adm2}(\iota,e,N|\operatorname{disc}K|,\rho),
                                                               \tag{RNC.35}
\]
all its original trace gaps are nonzero and (RNC.30) holds.
\begin{proof}
RNC11 constructs the finite catalogue with exactly
(RNC.35).  RNC13 constructs the integer solely from that
catalogue and the given indexed orbit.  RNC15 proves the
required size bound with constants chosen first.
RNC14 proves the original nonvanishing conjunction and
product inequality for every later coefficient datum
and family, at geometric $\sigma^{-1}$ as required.
The rank-two discriminant factor occurs only in its
original $M=N|\operatorname{disc}K|$; the rank-three
level is still $N$.  The auxiliary automorphic levels
used for counting have not altered these predicates.
These are all the clauses and the literal quantifier
order of the original $\mathrm{CNR3}$.
\end{proof}
% END RNC16

The proof leaves the original named classical and period
foundations in place, with the bounded applications above
checked explicitly.  It supplies the complete congruence
integer producer.  It does not supply the one-character
length theorem: the subsequent $\mathrm{PropGK}$ consumer
still needs that distinct arithmetic producer.
